\section*{الفصل \ref{ch:b2:alpine-orogen} --- تشريح حزام جبلي: جبال الألب}

\begin{solution}{exo:b2:alpine-orogen:1}
(أ) $e = 1 - 27.0/36.0 = 0.25$، أي \qty{25}{\percent}؛ $S = 36.0 - 27.0 =
\qty{9.0}{km}$. (ب) يجب أن تكون الطبقة قد احتفظت بطولها (دون تمدّد أو ترقّق على امتدادها)،
وأن يكون المقطع موازيًا لاتجاه النقل، وأن يحدّد خطّا التثبيت مواضع
لم تنزلق فيها الطبقات بعضها على بعض.
\end{solution}

\begin{solution}{exo:b2:alpine-orogen:2}
(أ) \omterm{def:b2:alpine-orogen:klippe}{كليب}: النيس بقية معزولة من غطاء زاحف مرتكزة على صخورها \omterm{def:b2:alpine-orogen:nappes}{الأصيلة}؛
وقد أزال الحتّ بقية الصفيحة من حولها كلها. (ب) \omterm{def:b2:alpine-orogen:klippe}{نافذة تكتونية}: قطع الحتّ
غطاء النيس الزاحف وكشف الصخور الأحدث تحت صدعه الدفعي القاعدي، الذي
ينحدر لذلك بعيدًا عن النافذة من كل الجهات.
\end{solution}

\begin{solution}{exo:b2:alpine-orogen:3}
(أ) المركّب فوق المافي (الوشاح العلوي)؛ والمركّب الغابروي (القشرة المحيطية السفلى)؛
ومجمع القواطع المصفوفة (القشرة الوسطى)؛ والمركّب البركاني، الوسادي (القشرة العليا).
(ب) مجمع القواطع المصفوفة. (ج) معظم قشرة المحيطات الواسعة يستهلكه الاندساس؛
والأوفيوليتات التي تبقى على اليابسة تكوّنت في الغالب في أحواض محيطية صغيرة أو بين أقواس
جزرية، والتيثيسية منها حيث رقّت قارة حتى انفتح محيط صغير.
\end{solution}

\begin{solution}{exo:b2:alpine-orogen:4}
(أ) غطاءا مونتي روزا وسان برنار الكبير: الجهة الأوروبية (البريانسوني
والهامش السلبي الأوروبي)؛ وزيرمات--ساس ونطاق كومبان: محيط بييمونت--ليغوريا؛ ومنظومة دان
بلانش: هامش أبوليا. (ب) أبوليا، التي تنتهي في القمّة. (ج) الصفيحة العليا
لا تندسّ، والشرائح المكشوطة من الصفيحة السفلى تتكدّس تحتها بالترتيب الذي
وصلت به الصفيحة السفلى: المحيط البعيد أولًا، والهامش القريب أخيرًا وأدنى. والطيّ اللاحق للركام،
أو الصدوع الدفعية الخارجة عن التسلسل، يمكن أن تُخلّ بالترتيب.
\end{solution}

\begin{solution}{exo:b2:alpine-orogen:5}
(أ) $2\pi G\,\Delta\rho = 2\pi\times6.674\times10^{-11}\times500 =
\qty{2.097e-7}{s^{-2}}$؛ $h = 1.70\times10^{-3}/2.097\times10^{-7} = \qty{8.1}{km}$.
(ب) $2.097\times10^{-7}\times25\,000 = \qty{5.24e-3}{m.s^{-2}}$، أي نحو
\qty{-524}{mGal}. (ج) الجذر ليس صفيحة لانهائية: بل جسم على عمق بضع عشرات من الكيلومترات
لا يزيد عرضه كثيرًا على عرض الحزام، فيجذب نقطة على السطح أقل مما تجذبها
صفيحة بالسمك نفسه ممتدة إلى ما لا نهاية؛ والأجسام الكثيفة كنطاق إيفريا
تضيف شذوذات موجبة، وقد يكون تباين الكثافة في العمق أصغر مما افتُرض.
\end{solution}

\begin{solution}{exo:b2:alpine-orogen:6}
(أ) $164 - 44 = \qty{120}{Myr}$. (ب) $z = P/(\bar\rho g)$:
$2.8\times10^9/(3000\times9.8) = \qty{95.2}{km}$؛ $2.8\times10^9/(3300\times9.8) =
\qty{86.6}{km}$. (ج) عبارة «أكثر من \qty{90}{km}» تقع داخل هذا المدى: فالعمق يتوقّف على
الكثافة المتوسطة المختارة للحمل العلوي بقدر ما يتوقّف على الضغط.
\end{solution}

\begin{solution}{exo:b2:alpine-orogen:7}
(أ) الجناح الذي طوله $L$ ومائل بالزاوية $\theta$ يمتدّ أفقيًا على $L\cos\theta$؛ ولكل الأجنحة
الميل نفسه، فيكون $l = l_0\cos\theta$ و$e = 1 - l/l_0 = 1 - \cos\theta$.
(ب) $\cos\theta = 0.70$: $\theta = \qty{45.6}{\degree}$. (ج) $1 - \cos\qty{40}{\degree} =
\qty{23.4}{\percent}$، كما في شكل الطيّات المتعرّجة؛ $1 - \cos\qty{60}{\degree}
= \qty{50}{\percent}$.
\end{solution}

\begin{solution}{exo:b2:alpine-orogen:8}
(أ) $l_0 = 118/2.0 = \qty{59.0}{km}$؛ $S = 59.0 - 42 = \qty{17.0}{km}$؛ $e = 17.0/59.0 =
\qty{28.8}{\percent}$. (ب) $A_e = 118 - 42\times2.0 = \qty{34.0}{km^2}$؛ $d = 34.0/17.0 =
\qty{2.0}{km}$. (ج) نعم: يقع \omterm{def:b2:alpine-orogen:decollement}{سطح الانفصال} عند قاعدة الغطاء الذي سمكه \qty{2.0}{km}، تحت
قمّته؛ فالغطاء منفصل على قاعدته.
\end{solution}

\begin{solution}{exo:b2:alpine-orogen:9}
(أ) جذر إيري هو $4.5h$: عند خط العرض \qty{46.5}{\degree} شمالًا، $4.5\times2.74 = \qty{12.3}{km}$، أي
موهو عند \qty{42.3}{km}؛ وعند \qty{45.55}{\degree} شمالًا، $4.5\times0.65 = \qty{2.9}{km}$، أي موهو
عند \qty{32.9}{km}. (ب) تحت الأرض المرتفعة يكون الموهو أقل عمقًا بمقدار \qty{6.0}{km} مما يتنبّأ به
إيري؛ وتحت حافة سهل البو أعمق بمقدار \qty{28.3}{km}. (ج) الجذر غير
معوَّض محليًا: فهو مزاح نحو الجنوب وأكبر مما تتطلّبه التضاريس.
فالصفيحة تنثني ككلّ تحت الأحمال التي عليها (الإسناد الإقليمي في فصل من مجلد السنة
الثالثة)، وأحد تلك الأحمال، جسم إيفريا، ليس تضاريس أصلًا.
\end{solution}

\begin{solution}{exo:b2:alpine-orogen:10}
(أ) $5.1/11.3 = 0.45$: أقل من النصف. (ب) $5.1 + 6.2 = \qty{11.3}{km}$. تقع قمّة
الجوراسي فوق \omterm{def:b2:alpine-orogen:decollement}{سطح الانفصال} العلوي؛ والقدر \qty{6.2}{km} المنقول على امتداد ذلك السطح
مرّ تحت طيّات المقطع واستُنفد في بنى أبعد شمالًا، فلم تتقصّر
الطبقة العليا بذلك المقدار داخل المقطع قطّ. (ج) الطبقة الدليلية الواحدة لا تسجّل
إلا تقصّر الطبقات التي تنتمي إليها؛ وكل سطح انفصال فوقها أو تحتها يمكن أن
يُدخل إزاحة إلى المقطع أو يُخرجها منه. وازِن عدة طبقات، وتحقّق بالمساحات.
\end{solution}

\begin{solution}{exo:b2:alpine-orogen:11}
(أ) لإزاحة قدرها \qty{100}{km}، $\sqrt{200^2 + 100^2} = \qty{223.6}{km}$؛ ولإزاحة
\qty{50}{km}، $\sqrt{200^2 + 50^2} = \qty{206.2}{km}$؛ ولإزاحة \qty{150}{km}،
$\sqrt{200^2 + 150^2} = \qty{250.0}{km}$. (ب) $\arctan(100/200) = \qty{26.6}{\degree}$ عن
الاتجاه شمال--جنوب، وهو خط تقارب قريب من الاتجاه \babelsublr{N\qty{150}{\degree}E} الذي ذكره روزنبرغ
ومعاونوه. (ج) الإزاحة غير مؤكدة بمقدار $\pm\qty{50}{km}$، ما يحرّك التقارب الكلي
بنحو \qty{20}{km} في كل اتجاه؛ والرقم الثالث في النتيجة أدنى بكثير من
عدم اليقين هذا.
\end{solution}

\begin{solution}{exo:b2:alpine-orogen:12}
(أ) من $210/12.9 = \qty{16.3}{Myr}$ إلى $210/8.5 = \qty{24.7}{Myr}$ (قسمة \unit{km} على
\unit{mm.yr^{-1}} تعطي ملايين السنين). (ب) أكثر من $4\times60 = \qty{240}{km}$؛
وعلى الأقل $240 - 60 = \qty{180}{km}$ من الحوض هي الآن تحت الصدوع الدفعية الألبية أو قد
حُتّت. (ج) يتكوّن الحوض أمام الموشور الدفعي ويهاجر معه: فجزؤه
الداخلي يُدسر فوقه دائمًا بالصدوع الدفعية التي تُثقله، وأقدم ملئه يعلق
في الصدوع الدفعية، كما هي حال \omterm{def:b2:alpine-orogen:foreland}{المولاس} تحت الألبي.
\end{solution}

\begin{solution}{pb:b2:alpine-orogen:1}
\noindent\textbf{الجزء الأول.}

\textbf{1.} $3.0/\cos\qty{35}{\degree} = \qty{3.66}{km}$؛ $2.5/\cos\qty{50}{\degree} =
\qty{3.89}{km}$؛ $4.0/\cos\qty{40}{\degree} = \qty{5.22}{km}$.

\textbf{2.} $8.00 + 2\times3.66 + 5.00 + 2\times3.89 + 6.0/\cos\qty{10}{\degree} +
2\times5.22 + 10.0/\cos\qty{5}{\degree} = 8.00 + 7.32 + 5.00 + 7.78 + 6.09 + 10.44 + 10.04 =
\qty{54.67}{km}$؛ ومع الامتداد المكرّر، $l_0 = \qty{60.7}{km}$.

\textbf{3.} الطبقة تحت الصفيحة الدفعية هي الطبقة نفسها، نقلها الصدع الدفعي تحت
الصفيحة: فقبل الدسر كانت على خط واحد مع بقيتها. وإغفالها يُرجع
الطبقة أقصر مما ينبغي.

\textbf{4.} $e = 1 - 48.0/60.7 = 0.209$، أي \qty{20.9}{\percent}؛ $S = 60.7 - 48.0 =
\qty{12.7}{km}$.

\textbf{5.} $60.7 - 60.2 = \qty{0.5}{km}$، $0.5/60.7 = \qty{0.8}{\percent}$: ضمن
دقة القياس؛ فالمقطع مقبول.

\textbf{6.} $60.7 - 54.0 = \qty{6.7}{km}$، أي \qty{11}{\percent}: لا يمكن لطبقتي صفيحة واحدة
أن تُرجعا إلى أطوال مختلفة إلى هذا الحدّ. فلن يكون المقطع مقبولًا كما رُسم: بنية
مفقودة أو مرسومة خطأً، أو خط تثبيت في المكان الخطأ، أو حركة خارج مستوى
المقطع، أو سطح انفصال بين الطبقتين.

\noindent\textbf{الجزء الثاني.}

\textbf{7.} $146/2.4 = \qty{60.8}{km}$.

\textbf{8.} يختلف عن \qty{60.7}{km} بمقدار \qty{0.1}{km}: فطول الخط والمساحة يتوافقان.

\textbf{9.} $A_e = 146 - 48.0\times2.4 = 146 - 115.2 = \qty{30.8}{km^2}$.

\textbf{10.} $d = 30.8/12.7 = \qty{2.43}{km}$، أي يساوي سمك الغطاء \qty{2.4}{km} في
حدود دقة البيانات: \omterm{def:b2:alpine-orogen:decollement}{فسطح الانفصال} هو قاعدة الغطاء.

\textbf{11.} $160/2.4 = \qty{66.7}{km}$، أي أكثر من طول الخط بمقدار \qty{6.0}{km}. فثخن
المتبخّرات عند قاعدة الغطاء، وهي تتدفّق إلى نوى الطيّات المحدّبة، يضيف مساحة
دون أن يضيف طولًا؛ فيفشل عندئذ توازن المساحة بسمك ثابت.

\textbf{12.} تحرّك الغطاء على سطح انفصال عند قاعدته ولم تشارك القاعدة.
وتحت الحزام ينبغي أن تمتدّ قمّة القاعدة مستوية أو منحدرة بلطف
وانتظام، دون أن تقطعها الصدوع الدفعية التي فوقها.

\noindent\textbf{الجزء الثالث.}

\textbf{13.} $22 + 18 + 25 = \qty{65}{km}$؛ $S = 65 - 15 = \qty{50}{km}$.

\textbf{14.} $e = 50/65 = \qty{76.9}{\percent}$، مقابل \qty{20.9}{\percent} في
حزام المقدّمة: فالأغطية أشدّ تقصّرًا بكثير.

\textbf{15.} $30 - 22 = \qty{8}{km}$ ($e = \qty{26.7}{\percent}$). فالقاعدة نفسها
متقصّرة: \omterm{def:b2:alpine-orogen:skin}{تكتونية سميكة الجلد}.

\textbf{16.} حزام المقدّمة $12.7/(11 - 4) = \qty{1.81}{km/Myr}$، أي
\qty{1.81}{mm} في السنة؛ والأغطية $50/(30 - 20) = \qty{5.0}{mm}$ في السنة.

\textbf{17.} تشوّهت الأغطية (من 30 إلى \qty{20}{Ma}) قبل حزام المقدّمة (من 11 إلى
\qty{4}{Ma}): فقد هاجر التشوّه نحو الخارج، نحو المقدّمة. وهذا ترتيب
الألب: دسر غلاروس في الأوليغوسين والميوسين، والكتل الخارجية في الميوسين،
والجورا في الميوسين المتأخر والبليوسين.

\textbf{18.} $12.7 + 50 + 8 = \qty{70.7}{km}$.

\noindent\textbf{الجزء الرابع.}

\textbf{19.} $\qty{0.5}{cm.yr^{-1}}\times\qty{40}{Myr} = 0.005\times4\times10^7 =
\qty{2.0e5}{m} = \qty{200}{km}$.

\textbf{20.} $200 + 18 = \qty{218}{km}$.

\textbf{21.} لا يُحسب إلا ما حُفظ أو استُنتج: فأجزاء من المحيط ربما
اندسّت قبل \qty{85}{Ma} أو أسرع من المعدّل المتوسط، وشرائح أوفيوليتية ربما
حُتّت أو نُقلت خارج المقطع، والتقارب المائل عن المقطع يُدخل
قاع محيط لا يراه المقطع.

\textbf{22.} النطاق الداخلي $60 - 20 = \qty{40}{km}$ ($e = \qty{66.7}{\percent}$)؛ والموشور الخلفي
$70 - 48 = \qty{22}{km}$ ($e = \qty{31.4}{\percent}$).

\textbf{23.} $12.7 + 50 + 8 + 40 + 22 = \qty{132.7}{km}$؛ ضمن مدى التقديرات
للألب الوسطى، نحو 125 إلى \qty{160}{km} من التقصّر التصادمي (قبل إضافة
الانزلاق المضربي).

\textbf{24.} العرض المُرجَع $153 + 132.7 = \qty{285.7}{km}$؛ والقشرة المُرجَعة
$285.7\times30 = \qty{8571}{km^2}$؛ والقشرة الحالية $153\times45 = \qty{6885}{km^2}$. المفقود:
$8571 - 6885 = \qty{1686}{km^2}$، أي \qty{19.7}{\percent} من القشرة المُرجَعة؛ وقد
حُتّ (وأُلقي في الأحواض الأمامية) أو اندسّ مع القشرة السفلى في
الوشاح. ولولا هذه الخسائر لكان سمك القشرة $8571/153 = \qty{56}{km}$ في المتوسط.
والتقصّر التصادمي الكلي عبر المقطع هو
\textbf{\qty[text-series-to-math=true]{133}{km}}.
\end{solution}
